% !TEX program = pdflatex
% !TEX encoding = UTF-8 Unicode

\documentclass[11pt, a4paper]{article}

% --- PACOTES PRINCIPAIS ---
\usepackage[utf8]{inputenc}
\usepackage[T1]{fontenc}
\usepackage{geometry}
\geometry{a4paper, top=2.5cm, bottom=2.5cm, left=2.0cm, right=2.0cm}
\usepackage{amsmath, amssymb, amsfonts}
\usepackage{booktabs}
\usepackage{tabularx}
\usepackage{graphicx}
\usepackage{xcolor}
\usepackage{enumitem}
\usepackage{tcolorbox}
\usepackage{fancyhdr}
\usepackage[colorlinks=true, linkcolor=ufopred, citecolor=deepnavy, urlcolor=ufopred]{hyperref}

% --- CORES INSTITUCIONAIS UFOP ---
\definecolor{ufopred}{RGB}{140, 29, 36}
\definecolor{deepnavy}{RGB}{20, 35, 60}
\definecolor{lightgray}{RGB}{245, 247, 250}
\definecolor{darkslate}{RGB}{30, 40, 50}

% --- CONFIGURAÇÃO DE CABEÇALHO E RODAPÉ ---
\pagestyle{fancy}
\fancyhf{}
\setlength{\headheight}{14pt}
\rhead{\textcolor{darkslate}{\small \textbf{Plataforma CCO/UFOP} --- Artigo Científico}}
\lhead{\textcolor{darkslate}{\small Modelagem Computacional Vascular}}
\cfoot{\thepage}

\tcbuselibrary{skins, breakable}

\begin{document}

% --- TÍTULO E AUTORES ---
\begin{center}
  {\LARGE \bfseries \textcolor{ufopred}{[Título do Artigo Científico em Português: Exemplo de Aplicação na Plataforma CCO/UFOP]}}\\[0.4cm]
  {\large \textbf{[Bernardo Sant' Anna Costa]}$^{1}$, \textbf{[Juliana Aparecida Borges]}$^{1}$, \textbf{[Sâmela Correa do Amaral]}$^{1}$,\textbf{[Rafael Alves Bonfim de Queiroz]}$^{1}$, }\\[0.2cm]
  {\small $^{1}$Departamento de Computação (DECOM) --- Universidade Federal de Ouro Preto (UFOP)\\[0.1cm]
  \texttt{\{bernardo.anna@aluno.ufop.edu.br, rafael.queiroz@ufop.edu.br\}}}\\[0.4cm]
\end{center}

% --- CAIXA DE RESUMO E PALAVRAS-CHAVE ---
\begin{tcolorbox}[colback=lightgray, colframe=deepnavy, arc=3pt, title=\textbf{RESUMO}, fonttitle=\bfseries]
\small
\textbf{Contexto:} Apresente sucintamente o problema fisiológico ou computacional abordado no contexto de modelagem de redes vasculares arteriais sintéticas ou híbridas. 
\textbf{Objetivo:} Descreva a principal contribuição do trabalho e o subprojeto da Plataforma CCO/UFOP no qual o estudo se insere (ex: CCO-GNN, CCO-Pulse, CCO-Retina).
\textbf{Metodologia:} Resuma os métodos computacionais, algoritmos de crescimento (Murray/Poiseuille), frameworks de inteligência artificial ou modelos de simulação hemodinâmica desenvolvidos e testados.
\textbf{Resultados:} Destaque os principais achados quantitativos, validações morfométricas, erros relativos e ganhos de desempenho computacional obtidos.
\textbf{Conclusão:} Sintetize o impacto científico e as potenciais aplicações médicas/tecnológicas dos resultados.

\vspace{0.2cm}
\textbf{Palavras-Chave:} Modelagem Vascular; Método CCO; Lei de Murray; Hemodinâmica Computacional; Plataforma CCO/UFOP.
\end{tcolorbox}

\vspace{0.4cm}

% --- SEÇÃO 1: INTRODUÇÃO ---
\section{Introdução}
As doenças cardiovasculares e cerebrovasculares permanecem como as principais causas de morbimortalidade no mundo. A simulação hemodinâmica não intrusiva e o planejamento cirúrgico dependem fundamentalmente da disponibilidade de representações geométricas precisas da árvore vascular humana. Contudo, imagens médicas distais (como angiotomografia) frequentemente apresentam limitações de resolução que impedem a reconstrução direta da microcirculação.

Neste cenário, a \textbf{Plataforma CCO/UFOP} desenvolve um ecossistema científico cumulativo baseado no algoritmo \emph{Constrained Constructive Optimization} (CCO), introduzido por Schreiner \& Buxbaum (1993). O método CCO permite a geração estocástica e construtiva de árvores vasculares sintéticas otimizadas para suprir um determinado domínio de perfusão sob restrições fisiológicas reais.

Este trabalho aborda o subprojeto \textbf{[Nome do Subprojeto, ex: CCO-Pulse / CCO-GNN]}, com o objetivo de [descrever a meta principal]. As principais contribuições deste artigo são:
\begin{itemize}[leftmargin=*]
  \item Proposição e implementação do algoritmo [descrever inovação];
  \item Validação morfométrica e hemodinâmica sob a Lei de Murray e equações de Navier-Stokes/Poiseuille;
  \item Integração ao repositório unificado no formato padronizado VTK.
\end{itemize}

% --- SEÇÃO 2: TRABALHOS RELACIONADOS ---
\section{Trabalhos Relacionados}
A literatura de geração automática de redes arteriais possui marcos importantes. O modelo clássico de CCO minimiza a potência vascular total $E$, equilibrando o trabalho de bombeamento hemodinâmico e o custo metabólico para manter o volume sanguíneo.

Trabalhos recentes estenderam o CCO para acomodar hemodinâmica pulsátil (Olufsen et al., 2000), calibração por otimização Bayesiana e modelos de inteligência artificial baseados em redes neurais em grafos (GNNs). A Tabela~\ref{tab:relacionados} resume a comparação entre o método proposto e abordagens da literatura.

\begin{table}[htbp]
  \centering
  \small
  \caption{Comparação qualitativa entre métodos da literatura e a proposta atual.}
  \label{tab:relacionados}
  \begin{tabularx}{\textwidth}{l X X X}
    \toprule
    \textbf{Abordagem} & \textbf{Modelo Hemodinâmico} & \textbf{Estrutura de Dados} & \textbf{Limitação / Lacuna} \\
    \midrule
    CCO Clássico & Estacionário (Poiseuille) & Árvore estrita (sem ciclos) & Omite fluxo pulsátil e colaterais \\
    Modelos 0D & Windkessel (Lumped) & Circuito elétrico equivalente & Ausência de geometria espacial \\
    \textbf{Proposto (CCO/UFOP)} & \textbf{[Estacionário / 1D Pulsátil]} & \textbf{[VTK Graph / Mesh]} & \textbf{[Contribuição deste trabalho]} \\
    \bottomrule
  \end{tabularx}
\end{table}

% --- SEÇÃO 3: FUNDAMENTAÇÃO TEÓRICA E METODOLOGIA ---
\section{Fundamentação Teórica e Metodologia}

\subsection{Modelo Fisiológico e Restrições de Bifurcação}
O escoamento sanguíneo nos segmentos vasculares cilíndricos é regido pela Lei de Poiseuille, onde a queda de pressão $\Delta P$ em um segmento de comprimento $L$, raio $r$ e vazão $Q$ é dada por:
\begin{equation}
  \Delta P = P_0 - P_1 = Q \cdot \left( \frac{8 \eta L}{\pi r^4} \right)
  \label{eq:poiseuille}
\end{equation}
sendo $\eta = 3.6 \times 10^{-3}\,\text{Pa}\cdot\text{s}$ a viscosidade dinâmica do sangue.

A relação de ramificação em cada bifurcação segue a Lei de Murray com expoente $\gamma$:
\begin{equation}
  r_0^\gamma = r_1^\gamma + r_2^\gamma \quad (2.55 \le \gamma \le 3.0)
  \label{eq:murray}
\end{equation}
onde $r_0$ é o raio do vaso pai, e $r_1, r_2$ são os raios dos vasos filhos.

\subsection{Algoritmo Desenvolvido e Arquitetura do Software}
Descreva nesta subseção o fluxo do algoritmo implementado no módulo \texttt{[cco\_xxx]}. Apresente o pseudocódigo ou o diagrama de blocos do processo computacional.

% --- SEÇÃO 4: RESULTADOS E DISCUSSÃO ---
\section{Resultados e Discussão}
Nesta seção, apresente os experimentos computacionais realizados. Inclua tabelas de métricas morfométricas e gráficos de convergência ou campos de pressão/fluxo.

\subsection{Análise Morfométrica e Validação}
A Tabela~\ref{tab:resultados} apresenta os resultados quantitativos obtidos para redes sintéticas com $N$ terminais.

\begin{table}[htbp]
  \centering
  \small
  \caption{Métricas morfométricas e desempenho computacional obtidos.}
  \label{tab:resultados}
  \begin{tabular}{ccccc}
    \toprule
    \textbf{Terminais ($N$)} & \textbf{Pressão de Perfusão (mmHg)} & \textbf{Vazão Total (mL/min)} & \textbf{Erro da Lei de Murray (\%)} & \textbf{Tempo de Execução (s)} \\
    \midrule
    50 & 100.0 & 120.5 & 0.12\% & 1.45 \\
    200 & 100.0 & 120.0 & 0.08\% & 8.32 \\
    1000 & 100.0 & 119.8 & 0.05\% & 45.10 \\
    \bottomrule
  \end{tabular}
\end{table}

\subsection{Discussão dos Achados}
Discuta a concordância dos resultados com dados fisiológicos conhecidos e analise o comportamento do algoritmo frente a limitações de memória ou tempo de processamento.

% --- SEÇÃO 5: CONCLUSÃO ---
\section{Conclusão}
Sintetize as principais conclusões do trabalho, ressaltando o cumprimento dos objetivos propostos. Indique os desdobramentos futuros para a Plataforma CCO/UFOP e sugestões de continuidade para novos alunos do grupo de pesquisa.

% --- AGRADECIMENTOS ---
\section*{Agradecimentos}
Os autores agradecem ao Departamento de Computação (DECOM) e à Universidade Federal de Ouro Preto (UFOP) pela infraestrutura computacional e suporte institucional. Este trabalho foi desenvolvido no âmbito da \textbf{Plataforma CCO/UFOP: Programa Integrado de Pesquisa em Modelagem Computacional de Redes Vasculares}.

% --- REFERÊNCIAS BIBLIOGRÁFICAS ---
\begin{thebibliography}{99}
  \bibitem{schreiner1993} SCHREINER, W.; BUXBAUM, P. Computer-optimization of vascular trees. \textbf{IEEE Transactions on Biomedical Engineering}, v. 40, n. 5, p. 482--491, 1993.
  \bibitem{murray1926} MURRAY, C. D. The physiological principle of minimum work: I. The vascular system and the cost of blood volume. \textbf{Proceedings of the National Academy of Sciences}, v. 12, n. 3, p. 207--214, 1926.
  \bibitem{olufsen2000} OLUFSEN, M. S. et al. Numerical simulation and experimental validation of blood flow in arteries with structured tree outflow conditions. \textbf{Annals of Biomedical Engineering}, v. 28, n. 11, p. 1281--1299, 2000.
  \bibitem{zamir1988} ZAMIR, M.; CHEE, H. Segment analysis of human coronary arteries. \textbf{Canadian Journal of Physiology and Pharmacology}, v. 66, n. 1, p. 155--159, 1988.
\end{thebibliography}

\end{document}
